\section*{Bab \ref{ch:b1:stress-strain-rheology} --- Tegangan, regangan dan reologi}

\begin{solution}{exo:b1:stress-strain-rheology:1}
(a) $\sigma_v = 2700\times9.81\times10\,000 = \qty{2.65e8}{Pa} = \qty{264.9}{MPa}$.
(b) $P = 1000\times9.81\times10\,000 = \qty{98.1}{MPa}$. (c) $\sigma_v' =
264.9 - 98.1 = \qty{166.8}{MPa}$, dan $\lambda = 98.1/264.9 = 0.370$, yang sama dengan
$1000/2700$.
\end{solution}

\begin{solution}{exo:b1:stress-strain-rheology:2}
(a) $e = (13-10)/10 = 0.30$ dan $S = 1.30$. (b) $e = (7-10)/10 = -0.30$, pemendekan
sebesar \qty{30}{\percent}, dan $S = 0.70$. (c) Sudut itu telah menyempit sebesar $\psi = \ang{30}$, dan $\gamma = \tan\ang{30} = 0.577$.
\end{solution}

\begin{solution}{exo:b1:stress-strain-rheology:3}
(a) Pusat \qty{70}{MPa}, jari-jari \qty{30}{MPa}. (b) Tegangan geser terbesar adalah jari-jari,
\qty{30}{MPa}, pada bidang-bidang yang membentuk \ang{45} dengan $\sigma_1$ dan $\sigma_3$, yang
membawa \omterm{def:b1:stress-strain-rheology:stress}{tegangan normal} pusat, \qty{70}{MPa}. (c)
$2\theta = \ang{60}$: $\sigma_n = 70 + 30\cos\ang{60} = \qty{85}{MPa}$ dan
$\tau = 30\sin\ang{60} = \qty{26.0}{MPa}$.
\end{solution}

\begin{solution}{exo:b1:stress-strain-rheology:4}
(a) $\ang{45} - \ang{30}/2 = \ang{30}$ terhadap $\sigma_1$, di kedua sisinya, sehingga kedua bidang
membentuk \ang{60} satu sama lain. (b) Dengan $c' = 20 + r$, syarat
\cref{thm:b1:stress-strain-rheology:failure-angle} menjadi
$r = C\cos\phi + (20 + r)\sin\phi$, sehingga
\[
  r = \frac{C\cos\phi + 20\sin\phi}{1 - \sin\phi} = \frac{8.66 + 10}{0.5}
    = \qty{37.3}{MPa},
\]
dan $\sigma_1' = 20 + 2r = \qty{94.6}{MPa}$.
\end{solution}

\begin{solution}{exo:b1:stress-strain-rheology:5}
(a) $c = 40$ dan $R = \sqrt{20^2 + 15^2} = \qty{25}{MPa}$: $\sigma_1 =
\qty{65}{MPa}$ dan $\sigma_3 = \qty{15}{MPa}$; $\tan2\theta_p = 30/40 = 0.75$,
$2\theta_p = \ang{36.87}$, $\theta_p = \ang{18.43}$. (b) Pada sudut ini
$\cos^2\theta = 0.9$, $\sin^2\theta = 0.1$ dan $\sin\theta\cos\theta = 0.3$:
$\sigma_n = 60\times0.9 + 2\times15\times0.3 + 20\times0.1 = 54 + 9 + 2 =
\qty{65}{MPa}$, dan $\tau = (20-60)\times0.3 + 15\times(0.9-0.1) = -12 + 12 =
0$. (c) $R = \qty{25}{MPa}$, pada bidang-bidang pada \ang{45} terhadap arah utama.
\end{solution}

\begin{solution}{exo:b1:stress-strain-rheology:6}
(a) $S_{1,2} = \sqrt{1 + 1} \pm 1$: $S_1 = 2.414$ dan $S_2 = 0.414$;
$\tan2\theta' = 2/2 = 1$, $\theta' = \ang{22.5}$. (b) $S_1S_2 =
(\sqrt2 + 1)(\sqrt2 - 1) = 1$, yang merupakan $\det F$: tidak ada perubahan luas. (c) Garis
sepanjang $(1, 1)$ menjadi $F(1, 1) = (1 + 2, 1) = (3, 1)$, pada
$\arctan(1/3) = \ang{18.4}$ dari arah geser, dengan \omterm{def:b1:stress-strain-rheology:strain}{peregangan}
$\sqrt{10}/\sqrt2 = 2.24$. Ia telah berputar melewati sumbu panjang elips,
pada \ang{22.5}, menuju bidang geser; \omterm{def:b1:stress-strain-rheology:strain}{peregangannya} sedikit kurang dari
$S_1$, seperti yang harus terjadi, karena tidak ada garis yang meregang lebih dari $S_1$.
\end{solution}

\begin{solution}{exo:b1:stress-strain-rheology:7}
(a) $S_1/S_2 = 4$ dan $S_1S_2 = 1$ memberi $S_1 = 2$ dan $S_2 = 0.5$; menurut
\cref{prop:b1:stress-strain-rheology:simple-shear}, $\gamma = S_1 - S_2 =
1.5$. (b) $\tan2\theta' = 2/1.5 = 1.333$, $2\theta' = \ang{53.1}$: sumbu panjang harus
terletak pada \ang{26.6} dari dinding. (c) $1.5\times200 = \qty{300}{m}$.
\end{solution}

\begin{solution}{exo:b1:stress-strain-rheology:8}
$\sigma_v = 2700\times9.81\times3000 = \qty{79.46}{MPa}$. (a) $P = \qty{29.43}{MPa}$,
$\sigma_n' = 70 - 29.43 = \qty{40.57}{MPa}$, dan Byerlee mengizinkan
$0.85\times40.57 = \qty{34.5}{MPa}$, lebih dari \qty{30}{MPa} yang diterapkan: sesar
bertahan. (b) $P = 0.6\times79.46 = \qty{47.68}{MPa}$, $\sigma_n' =
\qty{22.32}{MPa}$, dan gesekan hanya memberi \qty{19.0}{MPa}: sesar
tergelincir. (c) $30 = 0.85\,(70 - 79.46\,\lambda)$ memberi $\lambda = (70 -
35.29)/79.46 = 0.437$.
\end{solution}

\begin{solution}{exo:b1:stress-strain-rheology:9}
(a) Tegangan berubah sebagai $\dot\varepsilon^{1/4}$: $10^{1/4} = 1.78$. (b) Perbandingan kekuatannya adalah
\[
  \exp\Bigl[\frac{E}{nR}\Bigl(\frac{1}{673.15} - \frac{1}{773.15}\Bigr)\Bigr]
  = \exp\bigl(6705\times\num{1.921e-4}\bigr) = \exp(1.288) = 3.63 .
\]
(c) Perubahan laju sepuluh kali lipat mengubah kekuatan rayapan sebesar faktor 1.78, yang
dilakukan suhu hanya dalam $100\times\ln1.78/\ln3.63 \approx
\qty{45}{\degreeCelsius}$; pada \omterm{def:b1:heat-budget:geotherm}{geoterm} sekitar \qty{20}{\degreeCelsius}
per kilometer, itu menggeser transisi sekitar \qty{2}{km}. Perubahan \omterm{def:b1:heat-budget:geotherm}{geoterm} menggeser isoterm,
dan bersamanya transisi, jauh lebih banyak.
\end{solution}

\begin{solution}{exo:b1:stress-strain-rheology:10}
(a) Dengan $C = 0$, sesar tergelincir ketika $r = \mu c'/\sqrt{1+\mu^2}$
(\cref{thm:b1:stress-strain-rheology:failure-angle}). Di sini $c' = S =
\rho g z(1-\lambda)$, sehingga $\sigma_1 - \sigma_3 = 2r = 2\mu S/\sqrt{1+\mu^2}$.
(b) $S$ bertambah sebesar $2700\times9.81\times1000\times(1 - 1000/2700) =
\qty{16.68}{MPa}$ per kilometer, dan $2\times0.85/1.3124 = 1.295$: gradiennya
adalah \qty{21.6}{MPa} per kilometer, di antara \num{13.1} (sesar normal) dan
\qty{61.3}{MPa} per kilometer (sesar dorong). (c) $2\mu/(k+\mu) = 0.786$,
$2\mu/k = 1.295$ dan $2\mu/(k-\mu) = 3.68$, dengan $k = 1.3124$: \omterm{def:b1:stress-strain-rheology:principal}{tegangan
diferensial} saat tergelincir adalah 0.79, 1.30 dan 3.68 kali tegangan vertikal
efektif.
\end{solution}

\begin{solution}{exo:b1:stress-strain-rheology:11}
(a) $t = 60\times\num{e6}\times\num{3.156e7} = \qty{1.89e15}{s}$ dan
$2\sqrt{\kappa t} = 2\sqrt{\num{1.89e9}} = \qty{87.0}{km}$. (b)
$\operatorname{erf}x = 600/1350 = 0.444$ untuk $x = 0.417$: $z = 0.417\times87.0
= \qty{36.3}{km}$; $\operatorname{erf}x = 800/1350 = 0.593$ untuk $x = 0.586$:
$z = \qty{51.0}{km}$. (c) Gempa bumi seharusnya mencapai \qtyrange{36}{51}{km};
amplop menempatkan dasar lapisan rapuh pada \qty{38.2}{km}, di dalam rentang itu dan
dekat puncaknya, pada suhu \qty{641}{\degreeCelsius}, dalam rentang yang teramati.
\end{solution}

\begin{solution}{exo:b1:stress-strain-rheology:12}
(a) $r = 45$ dan $c' = \qty{65}{MPa}$; jarak dari pusat ke garis
$\tau = 0.85\sigma_n'$ adalah $0.85\times65/1.3124 = \qty{42.1}{MPa}$, kurang dari
jari-jari: lingkaran memotong garis, dan bidang-bidang pada busur di atasnya
tergelincir. (b) Dengan $2\theta$ sudut dari ujung $\sigma_1$, tergelincir memerlukan
$45\sin2\theta \ge 0.85\,(65 + 45\cos2\theta)$, yaitu
$45\sin2\theta - 38.25\cos2\theta \ge 55.25$, atau
$59.06\sin(2\theta - \ang{40.36}) \ge 55.25$, dengan $\ang{40.36} = \phi$.
Maka $\sin(2\theta - \phi) \ge 0.9355$, $2\theta - \phi$ antara \ang{69.31}
dan \ang{110.69}, $2\theta$ antara \ang{109.67} dan \ang{151.06}: normal-normalnya
terletak pada \ang{54.8} hingga \ang{75.5} dari $\sigma_1$, dan rekahan yang tergelincir
pada \ang{14.5} hingga \ang{35.2} dari $\sigma_1$. Titik tengah rentang itu,
\ang{24.8}, adalah $\ang{45} - \phi/2$, bidang
\cref{thm:b1:stress-strain-rheology:failure-angle}. (c) Hanya bidang itu yang tergelincir
ketika lingkaran hanya menyentuh garis: $r = \mu(20 + r)/k$, $r = 20\mu/(k -
\mu) = 17/0.4624 = \qty{36.8}{MPa}$, dan $\sigma_1' = 20 + 2r =
\qty{93.5}{MPa}$.
\end{solution}

\begin{solution}{pb:b1:stress-strain-rheology:1}
\textbf{1.} $\sigma_v = 2700\times9.81\times5000 = \qty{132.44}{MPa}$.

\textbf{2.} $P = 1000\times9.81\times5000 = \qty{49.05}{MPa}$, dan
$\lambda = 49.05/132.44 = 0.370$.

\textbf{3.} $\sigma_1 = \sigma_v = \qty{132.44}{MPa}$, vertikal, dan
$\sigma_3 = \qty{75.00}{MPa}$, horizontal timur--barat: \omterm{def:b1:layered-earth:layers}{kerak} sedang diregangkan timur--barat di bawah
beratnya sendiri, keadaan yang mendorong sesar normal, seperti sesar ini.

\textbf{4.} $\sigma_1 - \sigma_3 = \qty{57.44}{MPa}$; pusat
$(132.44 + 75.00)/2 = \qty{103.72}{MPa}$; jari-jari \qty{28.72}{MPa}.

\textbf{5.} $\sigma_1' = 132.44 - 49.05 = \qty{83.39}{MPa}$, $\sigma_3' =
75.00 - 49.05 = \qty{25.95}{MPa}$, pusat $103.72 - 49.05 = \qty{54.67}{MPa}$.

\textbf{6.} $\sigma_{xx} = \qty{75.00}{MPa}$, $\sigma_{yy} = \qty{132.44}{MPa}$
dan $\sigma_{xy} = 0$: sumbunya adalah arah utama, satu di antaranya vertikal, sehingga bidang-bidang
yang tegak lurus padanya tidak membawa tegangan geser.

\textbf{7.} Di sepanjang \omterm{def:b1:geological-maps:attitude}{kemiringannya} sesar berjalan sepanjang vektor satuan
$(\cos\ang{60}, -\sin\ang{60})$, ke timur dan ke bawah; normal ke atasnya adalah
$\vec n = (\sin\ang{60}, \cos\ang{60}) = (0.866, 0.500)$, pada \ang{30} dari
$x$. Kosinus sudutnya dengan $\sigma_1$, sepanjang $y$, adalah
0.500: $\theta = \ang{60}$.

\textbf{8.} Dengan \ang{30} dari $x$: $\sigma_n = 75.00\times0.75 + 132.44\times0.25
= 56.25 + 33.11 = \qty{89.36}{MPa}$, dan $\tau = (132.44 - 75.00)\times
\sin\ang{30}\cos\ang{30} = 57.44\times0.4330 = \qty{24.87}{MPa}$.

\textbf{9.} Pada $2\theta = \ang{120}$ dari ujung $\sigma_1$:
$\sigma_n = 103.72 + 28.72\cos\ang{120} = 103.72 - 14.36 = \qty{89.36}{MPa}$
dan $\tau = 28.72\sin\ang{120} = \qty{24.87}{MPa}$.

\textbf{10.} $\sigma_n' = 89.36 - 49.05 = \qty{40.31}{MPa}$.

\textbf{11.} $24.87/40.31 = 0.617$, di bawah 0.85: sesar bertahan.

\textbf{12.} $0.85\times40.31 = \qty{34.26}{MPa}$; tegangan geser kurang darinya sebesar
\qty{9.39}{MPa}.

\textbf{13.} Lingkaran memiliki pusat pada \qty{54.67}{MPa} dan jari-jari
\qty{28.72}{MPa}, dari \num{25.95} hingga \qty{83.39}{MPa}; garis Byerlee adalah
$\tau = 0.85\,\sigma_n'$ (semua \omterm{def:b1:stress-strain-rheology:stress}{tegangan normal} di sini di bawah
\qty{200}{MPa}); sesar berada pada titik $(40.31, 24.87)$, pada \ang{120} dari
ujung $\sigma_1$.

\textbf{14.} $0.85\times54.67/\sqrt{1 + 0.85^2} = 46.47/1.3124 =
\qty{35.41}{MPa}$, lebih besar dari jari-jari, \qty{28.72}{MPa}: seluruh lingkaran
berada di bawah garis, dan tidak ada bidang melalui batuan yang dapat tergelincir.

\textbf{15.} $\phi = \arctan0.85 = \ang{40.36}$: bidang pada
$\ang{45} - \phi/2 = \ang{24.8}$ dari $\sigma_1$ vertikal, dengan \omterm{def:b1:geological-maps:attitude}{kemiringan}
\ang{65.2}.

\textbf{16.} $S = \sigma_v' = \qty{83.39}{MPa}$ dan
$\sigma_1 - \sigma_3 = 2\times0.85\times83.39/(1.3124 + 0.85) =
\qty{65.55}{MPa}$; gradien memberi $13.11\times5 = \qty{65.55}{MPa}$.

\textbf{17.} $\sigma_3 = 132.44 - 65.55 = \qty{66.88}{MPa}$.

\textbf{18.} Sesar tergelincir ketika
$r\sin\ang{120} = 0.85\,(c' + r\cos\ang{120})$ dengan
$c' = (83.39 + \sigma_3')/2$ dan $r = (83.39 - \sigma_3')/2$; menyelesaikannya memberi
$\sigma_3' = \qty{17.18}{MPa}$, $\sigma_3 = \qty{66.23}{MPa}$. Ia lebih rendah
daripada untuk bidang terbaik: sesar tidak terletak di titik tempat lingkaran
pertama kali menyentuh garis, sehingga lingkaran harus tumbuh lebih jauh sebelum mencapai
titik sesar itu sendiri.

\textbf{19.} Setiap \omterm{def:b1:stress-strain-rheology:stress}{tegangan normal} efektif, termasuk $\sigma_1'$ dan $\sigma_3'$,
turun sebesar $\Delta P$, dan selisihnya tidak berubah; tegangan geser sama sekali tidak
bergantung pada tekanan pori.

\textbf{20.} Bidang terbaik tergelincir ketika jarak sama dengan jari-jari:
$0.85\,(54.67 - \Delta P)/1.3124 = 28.72$, sehingga
$\Delta P = 54.67 - 28.72\times1.3124/0.85 = 54.67 - 44.34 = \qty{10.33}{MPa}$.

\textbf{21.} Kolom air di dalam sumur beratnya sama seperti sebelumnya, sehingga
pompa harus menambahkan seluruh \qty{10.33}{MPa} di kepala sumur, sekitar 103 kali
tekanan atmosfer: berat kolom air
$\num{10.33e6}/(1000\times9.81) = \qty{1053}{m}$ tingginya.

\textbf{22.} Dengan $\mu = 0.6$ sesar hanya dapat menahan
$0.6\times40.31 = \qty{24.19}{MPa}$, lebih kecil dari \qty{24.87}{MPa} yang dibawanya:
ia sudah akan tergelincir. Karena ia tenang, gesekannya setidaknya $0.617$, atau
ia memiliki \omterm{def:b1:stress-strain-rheology:coulomb}{kohesi} tertentu.

\textbf{23.} Bidang terbaik tergelincir ketika tekanan pori mencapai
$49.05 + 10.33 = \qty{59.38}{MPa}$, yaitu $\lambda = 59.38/132.44 = 0.448$,
jauh di bawah nilai litostatik 1.

\textbf{24.} $\Delta P = \sigma_n' - \tau/0.85 = 40.309 - 24.870/0.85 = 40.309 -
29.259 = \qty{11.049}{MPa}$, lebih dari \qty{10.33}{MPa} yang akan menggerakkan
bidang orientasi terbaik. Hingga satu desimal, kenaikan tekanan pori yang membawa sesar
ke tergelincir adalah
\textbf{\qty[text-series-to-math=true]{11.0}{MPa}}.
\end{solution}
